% ---
% id: protocol-compiler/req:validator-refusals
% kind: req
% title: "Refusals of the validator's checks"
% summary: >-
%   A candidate declaration failing a check of the validator is refused, and
%   no declaration is emitted, with a reason naming the check and what that
%   check names. Nothing is repaired and no table is completed with a default
%   value.
% prominence: standard
% status: accepted
% uses: [def:validator]
% mentions: [thm:validator-soundness, foundations/req:consumer-decides-exactly]
% assumes: []
% refs: []
% provenance: newly-derived
% implements: certiqs_key.protocol_compiler.validation:validate
% lean:
% ---

A pass of every check of Definition~\ref{def:validator} is what establishes
the target conditions (Theorem~\ref{thm:validator-soundness}\,(i)), and the compiler, as a
\tref{consumer} of what it is handed, decides and never repairs
(\xref{foundations/req:consumer-decides-exactly}).

\begin{requirement}\label{req:validator-refusals}
The engine runs checks V1 to V8 of Definition~\ref{def:validator} on every
candidate declaration and refuses a candidate that fails one, emitting no
declaration, with a reason naming the failing check. It repairs nothing and completes no declared table with a default
value. Where a check has a reason of its own, the refusal carries that reason
and what the check names with it:
\begin{enumerate}[label=(\roman*),leftmargin=2.2em]
\item V1: \iface{sector_map_mismatch}, naming the first position at which the
  compared lists differ, and in the general format, for the block kind across
  the factors, the first factor at which the kind differs;
\item V2: \iface{alice_marginal_trace_rounds} where an addition of the
  marginal's trace rounds; in the general format
  \iface{alice_sector_marginal_zero}, naming the sector, where a block of the
  marginal is zero beside a measurement slot in \tref{ensemble form};
\item V3: \iface{alice_povm_summation_rounds}, in the general format carrying
  the Alice sector; \iface{bob_povm_summation_rounds}, carrying the sector and
  the outcome index at which the first inexact addition occurred, and in the
  general format the factor as well; in the general format
  \iface{alice_ensemble_marginal_mismatch}, carrying the sector and the first
  differing entry, \iface{alice_ensemble_sum_rounds}, carrying the sector and
  the index of the first inexact addition, and
  \iface{ensemble_form_in_untrusted_mode} for a slot in ensemble form beside a
  \tref{marginal slot} holding $\nomarg$;
\item V6: in the general format \iface{constraint_scope_not_a_factor},
  naming the entry;
\item V7: \iface{no_feasibility_witness}, naming, where a supplied witness was
  rejected on its map, the first position at which the two lists differ.
\end{enumerate}
\end{requirement}
